\chapter{Preliminaries}
\label{chap:preliminaries}
\label{sec:prelims}

For complex numbers $\alpha,\beta$ we write $\alpha\approx_\eps\beta$ if
$\abs{\alpha-\beta}\le\eps$; this relation satisfies the triangle inequality
with additive errors.

\section{Polynomials over finite fields}

There is a unique finite field $\F_q$ with $q$ elements for each prime power
$q$, and no other finite fields.

\begin{definition}[Polynomials of low individual degree]\label{def:polyfunc}
	\lean{MIPStarRE.LDT.Polynomial, MIPStarRE.LDT.PolynomialModel}
	\leanok
	Let $q$ be a prime power and $m,d\ge0$.  We write $\polyfunc mqd$ for the
	set of polynomials on $\F_q^m$ of individual degree at most $d$.
\end{definition}

\begin{lemma}[Schwartz--Zippel lemma]\label{lem:schwartz-zippel-total-degree}
	\lean{MIPStarRE.LDT.Preliminaries.schwartzZippel_totalDegree}
	\leanok
	Let $g,h\in\F_q[X_1,\ldots,X_m]$ be such that $g-h$ is nonzero of total
	degree at most $d$.  Then $\Prb_{x\sim\F_q^m}[g(x)=h(x)]\le d/q$.
\end{lemma}
\begin{proof}
	\leanok
	This is the Schwartz--Zippel lemma \cite{Sch80,Zip79}, available in
	Mathlib.
\end{proof}

\begin{corollary}[Schwartz--Zippel for individual degree]\label{cor:schwartz-zippel-individual}
	\lean{MIPStarRE.LDT.Preliminaries.polynomialAgreement_avg_le_mdq, MIPStarRE.LDT.Preliminaries.polynomialCollisionMass_le_mdq}
	\leanok
	\uses{def:polyfunc}
	If $g,h\in\polyfunc mqd$ are distinct, then
	$\Prb_{x\sim\F_q^m}[g(x)=h(x)]\le md/q$.  Consequently, for every
	sub-measurement $\{G_g\}$ on either side of a bipartite state,
	\[
		\sum_{g\ne h}\Prb_x[g(x)=h(x)]\,\bra\psi G_g\ot G'_h\ket\psi\le md/q.
	\]
\end{corollary}
\begin{proof}
	\leanok
	\uses{lem:schwartz-zippel-total-degree}
	The difference $g-h$ is nonzero of total degree at most $md$.  The second
	statement follows by summing the first against the positive weights
	$\bra\psi G_g\ot G'_h\ket\psi$, whose total is at most one.
\end{proof}

\section{Measurements}

All outcome sets are finite.

\begin{definition}[Measurements and sub-measurements]\label{def:measurements}
	\lean{MIPStarRE.LDT.SubMeas, MIPStarRE.LDT.Measurement, MIPStarRE.LDT.ProjSubMeas, MIPStarRE.LDT.ProjMeas}
	\leanok
	A sub-measurement on a Hilbert space $\calH$ is a family
	$A=\{A_a\}_{a\in\calA}$ of positive semidefinite operators with
	$\sum_aA_a\le I$.  It is projective if every $A_a$ is a projection, and it
	is a measurement if $\sum_aA_a=I$.  A family $A=\{A^x_a\}$ indexed by
	questions $x$ is called a (projective) sub-measurement or measurement if
	every $A^x$ is one.  We write $A^x=\sum_aA^x_a$ for the complete part.
\end{definition}

\begin{proposition}[Orthogonality of projective sub-measurements]\label{prop:projective-submeasurement-orthogonal}
	\lean{MIPStarRE.LDT.ProjSubMeas.outcome_orthogonal, MIPStarRE.LDT.ProjSubMeas.total_proj}
	\leanok
	\uses{def:measurements}
	If $P=\{P_a\}$ is a projective sub-measurement, then $P_aP_b=0$ for
	$a\ne b$, and $\sum_aP_a$ is a projection.
\end{proposition}
\begin{proof}
	\leanok
	With $S=\sum_bP_b\le I$, one has
	$0\le P_a(I-S)P_a=-\sum_{b\ne a}(P_bP_a)^\dagger(P_bP_a)$, so every
	$P_bP_a$ with $b\ne a$ vanishes.
\end{proof}

\begin{definition}[Polynomial measurements]\label{def:polymeasurements}
	\uses{def:measurements, def:polyfunc}
	We write $\polysub mqd$ for the sub-measurements with outcomes in
	$\polyfunc mqd$, and $\polymeas mqd$ for those among them that are
	measurements.
\end{definition}

\begin{definition}[Post-processing]\label{def:post-processing}
	\lean{MIPStarRE.LDT.postprocess}
	\leanok
	\uses{def:measurements}
	For a family $\{A_a\}_{a\in\calA}$ and a function $f\colon\calA\to\calB$,
	set $A_{[f(a)=b]}=\sum_{a:f(a)=b}A_a$.  For $G\in\polysub mqd$ and
	$u\in\F_q^m$, the evaluated sub-measurement is
	$\{G_{[g(u)=a]}\}_{a\in\F_q}$.
\end{definition}

\begin{proposition}[Post-processing preserves measurements]\label{prop:post-processing-preserves}
	\lean{MIPStarRE.LDT.postprocess_total}
	\leanok
	\uses{def:post-processing}
	$\sum_bA_{[f(a)=b]}=\sum_aA_a$.  In particular post-processing maps
	sub-measurements to sub-measurements and measurements to measurements.
\end{proposition}
\begin{proof}
	\leanok
	The fibres of $f$ partition $\calA$.
\end{proof}

\begin{definition}[Completing a sub-measurement]\label{def:measurement-completion}
	\lean{MIPStarRE.LDT.Preliminaries.completeAtOutcome}
	\leanok
	\uses{def:measurements}
	The completion of a sub-measurement $A$ adds the incomplete part
	$I-A^x$ as a further outcome.  The formalization assigns the incomplete
	part to a distinguished outcome $a_0$, obtaining the measurement
	$A_a+\bone{a=a_0}(I-A^x)$.
\end{definition}

\section{Comparing measurements}
\label{sec:comparing-measurements}

\begin{definition}[Consistency]\label{def:simeq}
	\lean{MIPStarRE.LDT.ConsRel}
	\leanok
	\uses{def:measurements}
	Let $\ket\psi\in\calH_{\mathrm A}\ot\calH_{\mathrm B}$ be a state, let
	$A=\{A^x_a\}$ and $B=\{B^x_a\}$ be sub-measurements on
	$\calH_{\mathrm A}$ and $\calH_{\mathrm B}$, and let $\calD$ be a
	distribution on questions.  We write $A^x_a\ot I\simeq_\delta I\ot B^x_a$
	if
	\[
		\E_{x\sim\calD}\sum_{a\ne b}\bra\psi A^x_a\ot B^x_b\ket\psi\le\delta.
	\]
\end{definition}

\begin{proposition}[Consistency for measurements]\label{prop:simeq-for-measurements}
	\lean{MIPStarRE.LDT.qBipartiteConsDefect_of_measurements}
	\leanok
	\uses{def:simeq}
	If $A$ and $B$ are measurements, then $A^x_a\ot I\simeq_\delta I\ot B^x_a$
	if and only if
	$\E_x\sum_a\bra\psi A^x_a\ot B^x_a\ket\psi\ge1-\delta$.
\end{proposition}
\begin{proof}
	\leanok
	By completeness the off-diagonal mass equals one minus the diagonal mass.
\end{proof}

\begin{definition}[The state-dependent distance]\label{def:approx_delta}
	\lean{MIPStarRE.LDT.SDDRel, MIPStarRE.LDT.SDDOpRel, MIPStarRE.LDT.Preliminaries.BipartiteSDDRel}
	\leanok
	For families of matrices $A=\{A^x_a\}$ and $B=\{B^x_a\}$ on $\calH$ we
	write $A^x_a\approx_\delta B^x_a$ if
	$\E_{x\sim\calD}\sum_a\|(A^x_a-B^x_a)\ket\psi\|^2\le\delta$.
\end{definition}

All estimates below are Hilbert-space inequalities in the weighted direct sum
of the vectors indexed by questions and outcomes: Cauchy--Schwarz, the
triangle inequality, and the bound
$\|(C^x_{a,b}v^x_a)\|\le\|(v^x_a)\|$ whenever
$\sum_b(C^x_{a,b})^\dagger C^x_{a,b}\le I$.

\begin{proposition}[Transferring consistency along the state-dependent distance]\label{prop:triangle-sub}
	\lean{MIPStarRE.LDT.Preliminaries.triangleSub_right, MIPStarRE.LDT.Preliminaries.triangleSub_right_heterogeneous}
	\leanok
	\uses{def:simeq, def:approx_delta}
	Let $A$ and $B$ be measurements and $C$ a sub-measurement.  If
	$A^x_a\ot I\simeq_\delta I\ot C^x_a$ and $A^x_a\ot I\approx_\eps B^x_a\ot I$,
	then $B^x_a\ot I\simeq_{\delta+\sqrt\eps}I\ot C^x_a$.  The formalization
	states the tensor-reversed form, with the sub-measurement on the left
	factor and the two measurements on the right factor.
\end{proposition}
\begin{proof}
	\leanok
	The two inconsistencies differ by
	$\E_x\sum_a\bra\psi(A^x_a-B^x_a)\ot C^x_a\ket\psi$, whose modulus is at
	most $\sqrt\eps$ by direct-sum Cauchy--Schwarz and $\sum_a(C^x_a)^2\le I$.
\end{proof}

\begin{proposition}[Rounding inside a consistency relation]\label{prop:round-right-submeasurement}
	\lean{MIPStarRE.LDT.Test.point_consistency_of_right_rounding, MIPStarRE.LDT.Test.point_consistency_of_left_rounding}
	\leanok
	\uses{def:simeq, def:approx_delta, def:post-processing}
	Let $A$ be a measurement on $\calH_{\mathrm A}$, let $X$ be a
	sub-measurement and $P$ a projective sub-measurement on
	$\calH_{\mathrm B}$, and let $f_x$ be functions.  If
	$A^x_a\ot I\simeq_\delta I\ot X^x_{[f_x(b)=a]}$ and
	$I\ot X^x_b\approx_\eta I\ot P^x_b$, then
	\[
		A^x_a\ot I\simeq_{2(\delta+\eta)}I\ot P^x_{[f_x(b)=a]},
	\]
	and the analogous statement holds with the tensor factors exchanged.
\end{proposition}
\begin{proof}
	\leanok
	Since each $P^x_b$ is a projection, the inconsistency with $P$ is the
	squared direct-sum norm of $((A^x_a)^{1/2}\ot P^x_b)\psi$ over pairs with
	$a\ne f_x(b)$.  Compare it with the corresponding vectors for $X$: their
	difference has squared norm at most $\eta$, while the vectors for $X$ have
	squared norm at most $\delta$ because $(X^x_b)^2\le X^x_b$.  The
	formalization uses $(\sqrt\delta+\sqrt\eta)^2\le2(\delta+\eta)$.
\end{proof}

\begin{proposition}[Consistency implies closeness for measurements]\label{prop:simeq-to-approx}
	\lean{MIPStarRE.LDT.Preliminaries.simeqToApprox, MIPStarRE.LDT.Preliminaries.simeqToApprox_heterogeneous, MIPStarRE.LDT.Preliminaries.approxToSimeq_heterogeneous}
	\leanok
	\uses{def:simeq, def:approx_delta}
	If $A$ and $B$ are measurements with $A^x_a\ot I\simeq_\delta I\ot B^x_a$,
	then $A^x_a\ot I\approx_{2\delta}I\ot B^x_a$.  For projective
	measurements the converse holds as well.
\end{proposition}
\begin{proof}
	\leanok
	Expanding the squared distance and using $A_a^2\le A_a$, $B_a^2\le B_a$,
	and completeness bounds it by $2-2\E_x\sum_a\bra\psi A_a^x\ot B_a^x\ket\psi$;
	for projective measurements the first inequality is an equality.
\end{proof}

\begin{proposition}[Closeness of inner products]\label{prop:closeness-of-ip}
	\lean{MIPStarRE.LDT.Preliminaries.closenessOfIP, MIPStarRE.LDT.Preliminaries.closenessOfIPAdjoint}
	\leanok
	\uses{def:approx_delta}
	If $A^x_a\approx_\gamma B^x_a$ and
	$\sum_a(\sum_bC^x_{a,b})(\sum_bC^x_{a,b})^\dagger\le I$ for all $x$, then
	\[
		\E_x\sum_{a,b}\bra\psi C^x_{a,b}A^x_a\ket\psi
		\approx_{\sqrt\gamma}
		\E_x\sum_{a,b}\bra\psi C^x_{a,b}B^x_a\ket\psi,
	\]
	and the adjoint version holds under the adjoint normalization.
\end{proposition}
\begin{proof}
	\leanok
	The difference is a direct-sum inner product of a vector of norm at most
	one with a vector of norm at most $\sqrt\gamma$.
\end{proof}

\begin{proposition}[Multiplying a close family by contractions]\label{prop:cab-approx-delta}
	\lean{MIPStarRE.LDT.Preliminaries.cabApproxDelta}
	\leanok
	\uses{def:approx_delta}
	If $A^x_a\approx_\delta B^x_a$ and $\sum_b(C^x_{a,b})^\dagger C^x_{a,b}\le I$
	for all $x,a$, then $C^x_{a,b}A^x_a\approx_\delta C^x_{a,b}B^x_a$.
\end{proposition}
\begin{proof}
	\leanok
	Apply the contraction bound to $v^x_a=(A^x_a-B^x_a)\psi$.
\end{proof}

\begin{proposition}[Triangle inequality for the state-dependent distance]\label{prop:triangle-inequality-for-approx_delta}
	\lean{MIPStarRE.LDT.Preliminaries.sddOpRel_chain, MIPStarRE.LDT.Preliminaries.stateDependentDistanceRel_triangle_three}
	\leanok
	\uses{def:approx_delta}
	If $(A_i)^x_a\approx_{\delta_i}(A_{i+1})^x_a$ for $i=1,\ldots,k$, then
	$(A_1)^x_a\approx_{(\sqrt{\delta_1}+\cdots+\sqrt{\delta_k})^2}(A_{k+1})^x_a$.
\end{proposition}
\begin{proof}
	\leanok
	Apply the direct-sum triangle inequality to the telescoping sum.
\end{proof}

\begin{proposition}[Triangle inequality for consistency]\label{prop:simeq-triangle-inequality}
	\lean{MIPStarRE.LDT.Preliminaries.simeqTriangleInequality_heterogeneous}
	\leanok
	\uses{def:simeq, prop:simeq-to-approx, prop:triangle-sub}
	Consistency relations between measurements on a bipartite state satisfy
	a triangle inequality with the square-root loss of
	Proposition~\ref{prop:triangle-sub}.
\end{proposition}
\begin{proof}
	\leanok
	Convert one relation to a state-dependent distance by
	Proposition~\ref{prop:simeq-to-approx} and transfer the other along it by
	Proposition~\ref{prop:triangle-sub}.
\end{proof}

\begin{proposition}[Consistency and closeness for sub-measurements]\label{prop:cons-sub-meas}
	\lean{MIPStarRE.LDT.Preliminaries.consSubMeas, MIPStarRE.LDT.Preliminaries.ConsSubMeasStmt}
	\leanok
	\uses{def:simeq, def:approx_delta}
	Let $A$ be a sub-measurement and $B$ a measurement with
	$A^x_a\ot I\simeq_\gamma I\ot B^x_a$.  Then
	$A^x_a\ot I\approx_\gamma A^x_a\ot B^x_a\approx_\gamma A^x\ot B^x_a$, and
	hence $A^x_a\ot I\approx_{4\gamma}A^x\ot B^x_a$.
\end{proposition}
\begin{proof}
	\leanok
	The differences $A^x_a\ot(I-B^x_a)$ and $(A^x-A^x_a)\ot B^x_a$ are
	positive contractions $X$ with $X^2\le X$, whose total expectation is the
	off-diagonal mass.
\end{proof}

\begin{proposition}[Switching a sandwich]\label{prop:switch-sandwich}
	\lean{MIPStarRE.LDT.Preliminaries.switchSandwich, MIPStarRE.LDT.Preliminaries.SwitchSandwichStmt}
	\leanok
	\uses{def:approx_delta}
	If $A=\{A^x_a\}$ is a projective sub-measurement with
	$A^x_a\ot I\approx_\delta I\ot A^x_a$, then for every $0\le B\le I$,
	\[
		\E_x\sum_a\bra\psi A^x_aBA^x_a\ot I\ket\psi
		\approx_{2\sqrt\delta}
		\E_x\sum_a\bra\psi B\ot A^x_a\ket\psi
		\approx_{\sqrt\delta}
		\E_x\sum_a\bra\psi BA^x_a\ot I\ket\psi.
	\]
\end{proposition}
\begin{proof}
	\leanok
	Compare the quadratic forms of $B\ot I$ on the direct-sum vectors
	$(A^x_a\ot I)\psi$ and $(I\ot A^x_a)\psi$, and use that the complete parts
	satisfy the same mirror bound.
\end{proof}

\begin{definition}[Strong self-consistency]\label{def:strong-self-consistency}
	\lean{MIPStarRE.LDT.SSCRel, MIPStarRE.LDT.BipartiteSSCRel}
	\leanok
	\uses{def:measurements}
	Let $\ket\psi\in\calH\ot\calH$ be permutation invariant.  A sub-measurement
	$A$ is $\delta$-strongly self-consistent if
	\[
		\E_x\sum_a\bra\psi A^x_a\ot A^x_a\ket\psi\ge\bra\psi A\ot I\ket\psi-\delta.
	\]
\end{definition}

\begin{proposition}[Two notions of self-consistency]\label{prop:two-notions-of-self-consistency}
	\lean{MIPStarRE.LDT.Preliminaries.twoNotionsOfSelfConsistency, MIPStarRE.LDT.Preliminaries.otherTwoNotionsOfSelfConsistency}
	\leanok
	\uses{def:strong-self-consistency, def:approx_delta}
	If $A$ is $\delta$-strongly self-consistent, then
	$A^x_a\ot I\approx_{2\delta}I\ot A^x_a$; for projective $A$ the converse
	holds as well.
\end{proposition}
\begin{proof}
	\leanok
	By permutation invariance the squared distance equals
	$2\E_x\sum_a\bra\psi(A^x_a)^2\ot I-A^x_a\ot A^x_a\ket\psi$, with equality
	in $(A^x_a)^2\le A^x_a$ for projective $A$.
\end{proof}

\begin{proposition}[Self-consistency after evaluation]\label{prop:two-notions-of-self-consistency-after-evaluation}
	\lean{MIPStarRE.LDT.Preliminaries.twoNotionsOfSelfConsistencyAfterEvaluation}
	\leanok
	\uses{prop:two-notions-of-self-consistency, def:post-processing}
	If $A$ is $\delta$-strongly self-consistent, then for every function $f$,
	$A^x_{[f(a)=b]}\ot I\approx_{2\delta}I\ot A^x_{[f(a)=b]}$.
\end{proposition}
\begin{proof}
	\leanok
	Post-processing keeps the complete part and can only increase the diagonal
	mass, so the post-processed family is again $\delta$-strongly
	self-consistent.
\end{proof}
